% ===========================================================================
%  Appendix B -- figures for this appendix go in figures/appendixB/
% ===========================================================================
\chapter{Compiling the document}\label{app:compile}
\chapterfigs{appendixB}


This template must be compiled with \textbf{XeLaTeX} (not pdfLaTeX), because it uses OpenType fonts and Unicode. The output is a PDF. Besides XeLaTeX, two helper programs are run: \textbf{BibTeX} for the bibliography and \textbf{makeglossaries} for the List of Acronyms.

\section{On your computer.}
With a full TeX Live or MiKTeX installation, run in the folder of the thesis
\begin{verbatim}
latexmk
\end{verbatim}
It reads the file \texttt{latexmkrc}, uses XeLaTeX and runs BibTeX and \texttt{makeglossaries} as many times as needed. \texttt{latexmk -c} deletes the auxiliary files. Without \texttt{latexmk}, the full sequence is
\begin{verbatim}
xelatex main
bibtex main
makeglossaries main
xelatex main
xelatex main
\end{verbatim}
\texttt{makeglossaries} is a Perl script; on Windows without Perl use \texttt{makeglossaries-lite main} instead. In any case the package option \texttt{automake} (in \texttt{preamble.tex}) builds the list by itself during compilation, so it usually works even if you forget this step.

\section{On Overleaf.}
Upload the zip (\emph{New Project} $\to$ \emph{Upload Project}). Then open the project \textbf{Menu} (top left) $\to$ \textbf{Settings} and change \textbf{Compiler} to \textbf{XeLaTeX}; the default compiler (pdfLaTeX) fails with this template. Overleaf runs BibTeX and \texttt{makeglossaries} automatically. If the output looks outdated or strange, use \emph{Recompile from scratch} to clear old auxiliary files. Free accounts have a compilation time limit that a large thesis can exceed; the tips below reduce the time.

\section{Overleaf on your own computer (free).}
Overleaf is open source: the Community Edition can be installed and run locally, so you do not need a paid plan. The code is at \href{https://github.com/overleaf/overleaf}{github.com/overleaf/overleaf}; the usual way to run it is with Docker, following the instructions of the \href{https://github.com/overleaf/toolkit}{Overleaf Toolkit}. Once it is running you use it in the browser, as the online version, and the settings above are the same. Some features of the paid versions are not included.


\section{Tips to compile faster while writing.}
\begin{itemize}
  \item \textbf{Compile only the chapter you are editing.} In \texttt{main.tex} uncomment the line \verb|\includeonly{chapters/chapter2}|. The other chapters are skipped but keep their numbers, citations and references.
  \item \textbf{Use Overleaf's Draft mode} (\emph{Recompile} menu): it skips the images and is much faster.
\end{itemize}